% revision_methods.tex -- working record of the methods added or changed in
% the first revision (PNAS 2026-35068). Each section is written so that its
% text can move into main.tex (Methods) or into an SI Methods section.
%
% Build: make methods  (revision_methods.pdf)
%
% Conventions:
%   - \status{...} prints a blue note on what is done, running, or open.
%     Define \hidestatus before \begin{document} to drop all of them.
%   - Each section ends with a comment block "Code / results": the pipeline
%     rule and script (repository kr-colab/ecolocator_pipeline, branch
%     main, directory simulations/) and where the outputs are. These
%     comments do not print.
%   - Reviewer points are numbered as in pnas_rev1/review-responses.tex:
%     R1.5 = reviewer 1, point 5.

\documentclass[11pt]{article}
\usepackage[utf8]{inputenc}
\usepackage[margin=1in]{geometry}
\usepackage{amsmath}
\usepackage{booktabs}
\usepackage{xspace}
\usepackage{xcolor}
\usepackage[hidelinks]{hyperref}
\usepackage[capitalize]{cleveref}
\usepackage[style=authoryear,citestyle=authoryear,sorting=nyt,natbib=true]{biblatex}
\addbibresource{ELrefs.bib}

\newcommand{\dfir}{Douglas-fir\xspace}
\newcommand{\ecolocator}{\texttt{EcoLocator}\xspace}
\newcommand{\slim}{\texttt{SLiM}\xspace}
\newcommand{\points}[1]{\par\noindent\textit{Reviewer points: #1.}\par}
\ifdefined\hidestatus
  \newcommand{\status}[1]{}
\else
  \newcommand{\status}[1]{\par\noindent\textcolor{blue}{[\textbf{Status:} #1]}\par}
\fi

\title{Revision methods: new and changed analyses}
\author{Working document, PNAS 2026-35068 revision}
\date{\today}

\begin{document}
\maketitle
\tableofcontents

%======================================================================
\section{Scope}

This document collects the methods that the revision adds or changes.
Sections~\ref{sec:pipeline}--\ref{sec:evaluation} describe each method in
full. Section~\ref{sec:corrections} lists statements in the submitted
Methods that are wrong or do not match the analysis. Status notes (blue)
record what is done and what is running; they are not manuscript text.

%======================================================================
\section{Simulation pipeline and reproducibility}
\label{sec:pipeline}

All continuous-space simulations and their analyses run in one Snakemake
workflow \citep{molder_snakemake_2021}. A table of conditions gives the
parameters of each condition and its number of replicates; the workflow
makes the environment maps, runs \slim, adds neutral variation, fits
\ecolocator, runs the genotype--environment association (GEA) tests, and
scores every method against the true QTL. Software versions are fixed by a
lock file: \slim 4.2.2 \citep{haller2023slim4}, \texttt{pyslim} 1.0.4,
\texttt{msprime} 1.1.1 \citep{baumdicker2022efficient}, LEA 3.22.0
\citep{frichot_lea_2015}, \texttt{shap} 0.51, and \ecolocator at a fixed
commit.

\paragraph{Seeds.} Every random step uses a seed derived from the condition
name and the replicate number (the first 8 hexadecimal digits of the
SHA-256 hash of ``\textit{condition}:\textit{replicate}''). \slim, recapitation,
the neutral mutations, and every leave-one-out fold use it, so a rerun gives
the same output.

\paragraph{Environment maps.} The 110 maps of the submitted analysis are
regenerated exactly with NLMpy 0.1.4 \citep{etherington_nlmpy_2015}:
midpoint displacement on an $8\times8$ grid with Hurst exponent $H=0$, map
$N$ with seed $N+41$. (NLMpy 1.2 gives different maps from the same seeds.)
Each map set also gets a report with the spatial autocorrelation of its
maps (Moran's $I$).

\paragraph{Model parameters.} The landscape is $5\times5$ units with
carrying capacity $K=100$ per unit area, so it holds about 2,500
individuals at most; the realized population is 1,300 to 1,700 individuals
(\cref{sec:runlength}), and every individual is sampled. QTL mutations
arise only inside the QTL windows of each architecture, at rate
$10^{-8}$ per base per generation, with normally distributed effects (SD
0.05); with the 200 kb window of architecture 1 the trait-wide QTL mutation
rate is $U = 2\times10^{-3}$ per genome copy. Architectures 2 and 3 have
300 kb and 500 kb of QTL windows, so their $U$ is 1.5 and 2.5 times
larger. Recombination is $10^{-6}$ per base, in \slim and in
recapitation (ancestral $N_e = 2{,}000$); neutral mutations are added at
$10^{-8}$ per base.

\points{R1.5 (population size, mutation rate, carrying capacity), R1.8 (maps)}
\status{Done. The pipeline reproduces the submitted LOO predictions
byte-for-byte on the same CPU type.}

% Code / results:
%   simulations/ecosim (command), workflow/Snakefile, config/config.yaml,
%   config/map_sets.tsv, workflow/scripts/make_maps.py, map_report.py,
%   overlay.py, to_ecolocator.py. Seeds: rep_seed() in the Snakefile.
%   Map check (Moran's I): map_reports/ of any run.

%======================================================================
\section{Run length and steady state}
\label{sec:runlength}

The submitted simulations ran for 10,000 ticks (the model default). We
tested whether that is long enough for local adaptation and neutral
diversity to reach steady state, in two ways.

\paragraph{Trajectories.} \slim writes, every 100 ticks, the population
size, the mean fitness from selection, the mean density regulation, the
correlation between phenotype $z$ and local optimum $\theta$, the mean
$|z-\theta|$, the phenotypic variance, the number of QTL mutations, and the
mean age. Writing the log draws no random numbers. For each replicate and
quantity, we average the log over blocks of 1,000 ticks (or $1/20$ of the
run, if shorter), take the plateau as the mean of the last 20\% of the run,
and the noise as the SD of the block means around a straight line fitted
to the second half. The time to plateau is the start of the first block
from which the block means stay within plateau $\pm 3\times$ noise for at
least 10\% of the run. A quantity whose change from the third to the last
quarter of the run exceeds 3 noise units has no plateau yet.

\paragraph{Coalescence.} Neutral diversity is added after \slim, so the
log cannot show whether it has reached equilibrium. For this, \slim writes
the tree sequence at fixed intervals during one long run (writing a
snapshot draws no random numbers; the \slim run is unchanged). For each
snapshot we compute the fraction of the genome whose tree has a single root
within the run. Below 1, part of the genealogy lies before the start of the
run, where recapitation models it as one well-mixed population without
space. Once the genome has coalesced within the run, the neutral genealogy
no longer depends on that assumption; this is the standard burn-in check
for tree-sequence simulations. We also compute, after recapitation and
without mutations, the expected neutral diversity and the slope of expected
divergence between spatial bins on log distance (isolation by distance).

\paragraph{Results.} In 10 runs of 500,000 ticks per dispersal level on the
$5\times5$ landscape ($\sigma_K=0.5$, architecture 1), 99\% of the genome
had coalesced within the run after a median of 80,000 ticks (90th
percentile 128,000; maximum 470,000) at $\sigma_D=0.1$, and 60,000 ticks
(maximum 105,000) at $\sigma_D=0.3$. At 10,000 ticks only 7\% and 21\% of
the genome had coalesced. We therefore rerun the Figure~2 conditions for
200,000 ticks, which covers about 90\% of replicates at $\sigma_D=0.1$ and
all replicates at $\sigma_D=0.3$; the longer runs add about 2\% to the
compute, because the leave-one-out fits take about 98\% of it.

In 5 replicates each, run length changed locus discovery much more than
prediction. From 10,000 to 100,000 ticks, LOO $R^2$ of location and
environment changed by less than 0.015, but the AUC-PR of every method for
ranking QTL fell by roughly half (for example, LFMM2 with $K=5$ from 0.37 to
0.20 at $\sigma_D=0.1$, and SHAP from 0.37 to 0.22 at $\sigma_D=0.3$). The
number of segregating QTL with minor-allele frequency (MAF) $\geq 5\%$ rose
from 20 to 31 at $\sigma_D=0.1$.

\points{R1.5 (long-term maintenance of locally adaptive alleles)}
\status{Burn-in: $5\times5$ done; $10\times10$ (1,000,000 ticks) running.
Steady-state rerun of the 7 Figure~2 conditions, 100 replicates each, at
200,000 ticks: running (started 2026-11-11).}

% Code / results:
%   SLiM log: LOGFILE in workflow/slim/localadaptation_nonWF.slim.
%   Plateau rule: workflow/scripts/trajectories.py (rule trajectories).
%   Snapshots: SNAPSHOT_INTERVAL in the SLiM model; workflow/scripts/
%   snapshot_stats.py (rule snapshot_stats); config/conditions_burnin.tsv.
%   Burn-in results: talapas /projects/kernlab/adkern/ecosim_snapshot_trial/
%     simulations/results/conditions_burnin/
%   10k vs 100k: talapas /projects/kernlab/adkern/ecolocator_pipeline/
%     equil_compare/equilibrium_summary.tsv
%   Steady-state rerun: config/conditions_steady.tsv; talapas
%     /projects/kernlab/adkern/ecolocator_pipeline_main/simulations/results/
%     conditions_steady/

%======================================================================
\section{Landscape size}
\label{sec:landscape}

To test whether the $5\times5$ landscape is too small, we compared it with
a $10\times10$ landscape at the same density ($K=100$ per unit area, so
about four times as many individuals), at $\sigma_D=0.1$ and 0.3
($\sigma_K=0.5$, architecture 1; 5 replicates each; 100,000 ticks). The
environment map grows with the landscape so that grid points stay 0.71
units apart: $8\times8$ for $5\times5$ (the submitted maps) and
$15\times15$ for $10\times10$, both midpoint displacement with $H=0$. At
the same $H$, NLMpy's larger maps are smoother (Moran's $I$ about 0.49
against 0.37).

\status{Trajectories done (simulation only). The $10\times10$ burn-in runs
(coalescence) end about 2026-11-12. No LOO fits on $10\times10$ yet.}

% Code / results:
%   config/conditions_landscape.tsv (loo False); map sets mpd_h0_8x8 and
%   mpd_h0_15x15 in config/map_sets.tsv.

%======================================================================
\section{Polygenicity}
\label{sec:polygenicity}

To cover the regime of many small-effect alleles, we varied the trait-wide
QTL mutation rate $U$ (QTL mutation rate times the 200 kb QTL window of
architecture 1) over three orders of magnitude, from few large-effect to
many small-effect QTL, at $\sigma_D=0.3$ and $\sigma_K=0.5$. The SD of QTL
effects scales with $1/\sqrt{U}$, so the new genetic variance per
generation, $U\sigma_\alpha^2$, stays the same:

\begin{center}
\begin{tabular}{lccc}
\toprule
Condition & $U$ & QTL mutation rate & SD of QTL effects \\
\midrule
u2e-4 & $2\times10^{-4}$ & $10^{-9}$ & 0.158 \\
u2e-3 & $2\times10^{-3}$ & $10^{-8}$ (default) & 0.05 (default) \\
u2e-2 & $2\times10^{-2}$ & $10^{-7}$ & 0.0158 \\
u2e-1 & $2\times10^{-1}$ & $10^{-6}$ & 0.005 \\
\bottomrule
\end{tabular}
\end{center}

\noindent Each condition has 10 replicates of 200,000 ticks. We measure the
architecture of each replicate (\cref{sec:architecture}): the number of QTL
sites that carry local adaptation and their ages.

\points{R1.5 (critical mutation rate; many alleles of small effect), R2.5
(architectures not polygenic enough), R1.14, R1.15}
\status{A pilot at 10,000 ticks: the effective number of contributing
sites rose from 2.4 to 589 across the four rates; LOO $R^2$ did not change;
the SHAP ranking of QTL fell to random at high $U$; LFMM2 with $K=1$ held.
The 200,000-tick runs are running (started 2026-11-11).
Open: the critical mutation rate of Sakamoto et al.\ (2025, \textit{Genetics})
for this population structure, to place the four rates relative to it.}

% Code / results:
%   config/conditions_polygenic.tsv (columns m, qtl_sd, runtime).
%   Pilot (10k): talapas /projects/kernlab/adkern/ecolocator_pipeline/
%     polygenic_pilot/
%   Rerun: talapas .../ecolocator_pipeline_main/simulations/results/
%     conditions_polygenic/
%   TODO: add the Sakamoto et al. 2025 entry to ELrefs.bib
%     (Genetics, doi 10.1093/genetics/iyae165).

%======================================================================
\section{Leave-one-out cross-validation with per-fold SHAP}
\label{sec:loo}

Every individual is one fold. A model trained on all other individuals
predicts its location $(x, y)$ and environment, and gives the SHAP values
of its genotype, so every prediction and every SHAP value comes from a
model that did not see that individual. In each fold, 10\% of the training
individuals are held out for validation (early stopping after 100 epochs
without improvement, at most 5,000 epochs; the learning rate drops after
16 epochs without improvement; batch size 32). The network is the
\ecolocator default: batch normalization, 10 dense layers of 256 ELU
units with dropout 0.25 after the fifth, and two output heads; loss
weights $w_\text{loc}=w_\text{env}=0.5$. A fold's validation split and
network initialization depend only on the replicate seed and the fold
index, so the result does not depend on how folds are distributed across
processes.

SHAP values are expected gradients (\texttt{shap.GradientExplainer},
\citealp{lundberg_shap_2017}) of each output with respect to the held-out
genotype, with 100 training individuals of that fold as the background.
The SNP score is the mean $|\text{SHAP}|$ over all folds, for each output
($x$, $y$, environment). We rank SNPs by this raw score; no step uses the
QTL labels (\cref{sec:corrections}, item~3).

\points{R1.2 (what counts as an adaptive allele)}
\status{Done (package \texttt{ecolocator loo} with \texttt{--fold-start},
\texttt{--fold-stop}, \texttt{--shap-out}). A pilot showed that 5-fold
cross-fitted SHAP matches LOO SHAP, and that conditioning SHAP on genotype
PCs makes the QTL ranking worse.}

% Code / results:
%   Package: kr-colab/ecolocator, EcoLocator.fit_predict_loo (fold ranges,
%   per-fold SHAP of the held-out sample).
%   Pipeline: rule loo, workflow/scripts/loo_parallel.py, merge_loo.py,
%   merge_shap.py. Outputs per replicate: loo_predictions.tsv, metrics.tsv,
%   shap/mean_abs_shap.tsv.

%======================================================================
\section{LFMM2}
\label{sec:lfmm}

We ran LFMM2 \citep{caye_lfmm2_2019} with LEA 3.22.0
(\texttt{lfmm2} and \texttt{lfmm2.test}) on the individual genotypes
(0/1/2) against the local optimum, for $K = 1$--6, 8, and 10 latent
factors. The scree plot of the genotype PCA has no clear elbow in these
simulations, so we keep results for all $K$ and report $K=1$ and $K=5$.
Sites with a minor-allele count below 2 are not tested. $p$-values are
calibrated by the genomic inflation factor (\texttt{genomic.control =
TRUE}). For the overlap with SHAP, LFMM2 hits are sites with
Benjamini--Hochberg $q < 0.05$ \citep{benjamini_hochberg_1995}.

\points{R1.1, R1.9 (compare with LFMM2)}
\status{Done.}

% Code / results:
%   workflow/scripts/lfmm.R (rule lfmm). Outputs per replicate:
%   lfmm/pvalues.tsv, lfmm/summary.tsv (GIF per K), lfmm/eigenvalues.tsv.

%======================================================================
\section{Location-corrected GEA}
\label{sec:correctedgea}

The location-corrected GEA uses \ecolocator's predicted locations as the
correction for population structure. For every SNP $j$ we fit, by least
squares,
\begin{equation*}
g_j = b_0 + b\,e + f(\hat{x}, \hat{y}) + \varepsilon,
\end{equation*}
where $g_j$ is the 0/1/2 genotype, $e$ the environment, and $f$ a smooth
two-dimensional function (tensor-product cubic B-splines, up to 5 knots per
axis, fewer for small samples, at least 10 individuals per basis function)
of the location that \ecolocator predicted for the individual with a model
that did not see it (the LOO predictions, \cref{sec:loo}). The predicted
location stands for population structure; the true location would also
carry the environment and remove the signal. The $z$-scores
$\hat{b}/\mathrm{SE}(\hat{b})$ are calibrated by the genomic inflation factor,
as in LFMM2.

\status{Done. Over 873 replicates of the submitted-style runs, it ranked
QTL as well as LFMM2 with $K=5$ (mean difference in AUC-PR $-0.001$), and
better than LFMM2 with $K=1$ ($+0.028$; better in 69\% of replicates), SHAP
($+0.057$; 75\%), and the same test with true locations ($+0.030$; 78\%).
Candidate for the main locus-discovery result.}

% Code / results:
%   workflow/scripts/location_corrected_gea.py (rule corrected_gea).
%   Outputs per replicate: corrected_gea/results.tsv, summary.tsv.
%   873-replicate comparison: talapas /projects/kernlab/adkern/
%     ecolocator_pipeline/audit/corrected_gea.tsv

%======================================================================
\section{WZA and the top-candidate test}
\label{sec:wza}

We follow \citet{booker_wza_2024} and the documented use of their script
(\texttt{general\_WZA\_script.py}, TBooker/WZA on GitHub, commit
\texttt{8191f58}; the same file as in the submitted analysis).
\begin{enumerate}
\item For every SNP, Kendall's $\tau_b$ between genotype (0/1/2) and the
  environment of each individual, with its $p$-value. $\tau_b$ allows for
  the many ties of genotype data.
\item The script removes SNPs with MAF $\leq 0.05$ and turns each
  $p$-value into an empirical $p$-value: its rank among all SNPs divided by
  the number of SNPs.
\item Windows are 10 kb, the window size of Booker et al.'s simulations.
  The WZA combines the empirical $p$-values of a window into a weighted $Z$
  (Stouffer's $Z$ with weights $\bar{p}(1-\bar{p})$, where $\bar{p}$ is the
  allele frequency), and the script corrects its $p$-value for the number
  of SNPs per window.
\item The top-candidate test counts the SNPs in a window whose empirical
  $p$-value is in the top 1\% genome-wide, and tests for an excess with a
  binomial test whose expectation is 1\%.
\end{enumerate}

\noindent The empirical $p$-values are part of the method. The
top-candidate test also depends on them, because its binomial expectation
is the hit threshold itself. The submitted analysis passed parametric
$p$-values with the script's \texttt{--empiricalP} flag, which skips step
2, and used the top 5\% as hits (\cref{sec:corrections}, item~4).

Booker et al.\ analysed allele frequencies of sampled demes; here every
individual's genotype is correlated with its own environment. The
empirical ranking is the same in both cases.

\points{R1.9 (WZA was designed for allele frequencies)}
\status{Done. On one replicate of $\sigma_D=0.1$, the documented use raised
the window-level AUC-PR of WZA from 0.53 to 0.63 and of the top-candidate
test from 0.29 to 0.68. Open: an allele-frequency version (sampled demes)
to test Booker et al.'s design directly.}

% Code / results:
%   workflow/scripts/wza_input.py (Kendall tau-b, MAF, windows), rule wza.
%   Rule fetch_wza downloads Booker's script at the fixed commit (MD5
%   a2fa7df791f9b817779585ea440a9d74). Outputs: wza/input.csv,
%   wza/windows.csv.

%======================================================================
\section{Genetic architecture and true positives}
\label{sec:architecture}

Phenotypes are additive sums of QTL effects, so the covariance between
phenotype $z$ and local optimum $\theta$ over all individuals splits
exactly into one term per QTL site:
\begin{equation*}
\mathrm{Cov}(z,\theta) = \sum_i c_i, \qquad c_i = \mathrm{Cov}(g_i, \theta),
\end{equation*}
where $g_i$ is an individual's summed effect at site $i$ (both copies). We
call $c_i$ the contribution of site $i$ to local adaptation. A negative
$c_i$ marks a site that works against local adaptation. From the
contributions we report the number of segregating QTL sites, the effective
number of contributing sites $(\sum|c_i|)^2/\sum c_i^2$, the shares of the
largest 1 and 5 sites, and the contribution-weighted mean age.

We score every method against three kinds of true positive:
\begin{description}
\item[Any QTL.] The SNP is a QTL site; the window holds one.
\item[Major QTL.] The fewest QTL sites that together carry 80\% of
  $\sum_i |c_i|$. The sign is not used, because a GEA test detects
  association in either direction. This set is relative: every replicate
  has major QTL.
\item[Locally adapted (Booker et al.'s rule).] A window is locally adapted
  if the summed contribution of its QTL sites (the covariance between the
  window's additive effect and $\theta$) exceeds 0.005, the threshold of
  \citet{booker_wza_2024} for genes; a SNP, if its own $c_i$ does. Only
  positive contributions qualify, and the threshold is absolute, so a
  weakly adapted replicate can have none. Booker et al.\ computed the
  covariance from deme means; here it is from individuals, each with its
  own optimum.
\end{description}

\noindent In the 10,000- and 100,000-tick runs, the major QTL were 8 to 12
sites per replicate (15--20\% of QTL sites), and the threshold of 0.005
marked 5.6 to 6.8 windows per replicate, close to the 5.4 to 7.2 locally
adapted genes per replicate of \citet{booker_wza_2024}.

\points{R1.2, R1.16 (what counts as an adaptive allele)}
\status{Done.}

% Code / results:
%   workflow/scripts/architecture.py (rule architecture): architecture/
%   qtl_contributions.tsv, architecture/summary.tsv. The script checks that
%   the per-site effects add up to the phenotypes SLiM recorded.
%   Major QTL and Booker rule: major_qtl() and evaluate() in
%   workflow/scripts/compare_rep.py; settings major_qtl_share,
%   adapted_cov_threshold in config/config.yaml.

%======================================================================
\section{Evaluation of the SNP and window rankings}
\label{sec:evaluation}

\paragraph{Scores.} Every method gives each SNP a score in which larger
means stronger evidence and no score uses the QTL labels: mean
$|\text{SHAP}|$ for each output; $-\log_{10} p$ for LFMM2 (each $K$), the
location-corrected GEA, and Kendall's $\tau_b$; and for WZA and the
top-candidate test, the $-\log_{10} p$ of the SNP's window. SNPs with MAF
below 5\% are left out.

\paragraph{Levels.} We evaluate the rankings at two levels.
At the SNP level, every SNP is one unit.
At the window level, the 10 kb WZA windows are the units: WZA and the
top-candidate test score windows themselves, and every other method scores
a window by its strongest SNP, as in the comparison of
\citet{booker_wza_2024}. The SNP level favours single-SNP methods and the
window level favours window methods, so we report both.

\paragraph{Metrics.} For each replicate, level, kind of true positive
(\cref{sec:architecture}), and method we compute the area under the
precision--recall curve (AUC-PR; trapezoidal, as in the submitted
analysis). We summarize each condition by the mean, SD, and standard error
over replicates. For precision--recall and ROC curves of a condition, we
pool the replicates after turning each method's scores into percentile
ranks within each replicate, because raw scores (mean $|\text{SHAP}|$ above
all) are on a different scale in every replicate.

\points{R1.1, R1.9}
\status{Done. Preliminary (5 replicates per group; window level; major QTL
as true positives): at 100,000 ticks, Kendall's $\tau_b$, LFMM2, and the
location-corrected GEA reach AUC-PR 0.65--0.80; the top-candidate test
0.64--0.72; WZA 0.36--0.43; SHAP 0.22--0.50. Consistent with the WZA
preprint's own result that single-SNP tests win when allele frequencies
are well estimated and the environment is the selective agent.}

% Code / results:
%   workflow/scripts/compare_rep.py (rule compare): comparison/snps.tsv,
%   metrics.tsv, eval.tsv, windows.tsv, snp_scores.tsv.
%   workflow/scripts/summarize_comparison.py (rule comparison_all):
%   comparison_eval_summary.tsv (mean/SD/SE per condition, level, truth,
%   method) and the tidy tables for figures.
%   workflow/scripts/plot_comparison.py (rule comparison_curves): PR/ROC
%   figures per condition, comparison_curves.tsv, comparison_curves_auc.tsv.
%   Preliminary tables: poppy ~/ecolocator_pipeline/simulations/results/
%   equil_compare/

%======================================================================
\section{Corrections to the submitted Methods}
\label{sec:corrections}

\begin{enumerate}
\item \textbf{Spatial autocorrelation of the maps.} The Methods state that
  $H=0$ ``removes spatial autocorrelation''. The 110 maps have Moran's $I$
  of 0.36 on average (0.08--0.61). The environment is patchy, not
  uncorrelated with space. (R1.8)
\item \textbf{Run length.} The Methods do not state the run length. It was
  10,000 ticks, far from steady state (\cref{sec:runlength}).
\item \textbf{SHAP ``empirical FDR score''.} The score ranks SNPs by mean
  $|\text{SHAP}|$ and then assigns each rank the fraction of neutral SNPs
  above it, computed from the true labels. Ranking by this score uses the
  labels being evaluated. The revised analysis ranks by raw mean
  $|\text{SHAP}|$.
\item \textbf{WZA and top-candidate settings.} The Methods describe
  parametric $p$-values and a top-1\% threshold. The analysis used
  \texttt{--empiricalP} with parametric $p$-values (no empirical step) and
  a top-5\% threshold (\cref{sec:wza}).
\item \textbf{Loss weights.} The Methods give a default of 1.0 for both
  weights (true for the \ecolocator command line); the simulations used
  0.5 for both. The ratio is the same.
\item \textbf{Selection sweep values.} The Methods give $\sigma_K = 0.3,
  0.5, 0.7$; Fig.~S4 gives 0.5, 0.7, 0.9. The simulations used 0.3, 0.5,
  0.7.
\item \textbf{Booker et al.'s comparison.} The Methods say that Booker et
  al.\ found WZA and top-candidate as good as or better than Kendall's
  $\tau$, BayPass, LFMM, and RDA. Their preprint compares Kendall's $\tau$
  and BayPass only, and its Fig.~S12 shows single-SNP tests ahead when
  allele frequencies are known well.
\item \textbf{Comparison level.} The submitted comparison mapped window
  $p$-values onto SNPs and scored at the SNP level only, which puts window
  methods at a disadvantage (\cref{sec:evaluation}).
\end{enumerate}

\status{Check item 6 against Jordan's run configuration, and item 7
against the published version of \citet{booker_wza_2024}, before changing
the text.}

\printbibliography

\end{document}
